% problem-000428 1 合成氨原料气制备
\section*{1. 合成氨原料气制备}
\textit{下列为分问。}

\subsection*{(1)}
$\mathrm { C H _ { 4 } ( g ) + H _ { 2 } O ( g )  C O ( g ) + 3 H _ { 2 } ( g ) }$

\subsection*{(2)}
$2 \mathrm { C H } _ { 4 } ( \mathrm { g } ) + \mathrm { O } _ { 2 } ( \mathrm { g } ) \to 2 \mathrm { C O } ( \mathrm { g } ) + 4 \mathrm { H } _ { 2 } ( \mathrm { g } )$

$
(3) \quad \mathrm{CO} (\mathrm{g}) + \mathrm{H} _ {2} \mathrm{O} (\mathrm{g}) \rightarrow \mathrm{H} _ {2} (\mathrm{g}) + \mathrm{CO} _ {2} (\mathrm{g})
$

假设反应产⽣的CO全部转化为 $\mathrm { C O _ { 2 } , \ C O _ { 2 } }$ 被碱液完全吸收，剩余的 $\mathrm { H _ { 2 } O }$ 通过冷凝⼲燥除去。进⼊合成氨反应塔的原料⽓为纯净的 $\mathrm { N } _ { 2 }$ 和 $\mathrm { H } _ { 2 \circ }$

\subsection*{1-1}
为使原料⽓中N_{2}和H_{2}的体积⽐为1: 3，推出起始⽓体中 $\mathrm { C H _ { 4 } }$ 和空⽓的⽐例。设空⽓中 $\mathrm { O _ { 2 } }$ 和 $\mathrm { N } _ { 2 }$ 的体积⽐为1:4，所有⽓体均按理想⽓体处理。

\subsection*{1-2}
计算反应(2)的反应热。已知：

$
\mathrm{C} (\mathrm{s}) + 2 \mathrm{H} _ {2} (\mathrm{g}) \rightarrow \mathrm{CH} _ {4} (\mathrm{g}) \quad \Delta H _ {4} = - 7 4. 8 \mathrm{kJ} \mathrm{mol} ^ {- 1} \tag {4}
$

$
\mathrm{C} (\mathrm{s}) + 1 / 2 \mathrm{O} _ {2} (\mathrm{g}) \rightarrow \mathrm{CO} (\mathrm{g}) \quad \Delta H _ {5} = - 1 1 0. 5 \mathrm{kJ} \mathrm{mol} ^ {- 1} \tag {5}
$
